\documentclass[12pt,a4paper]{article}

% --- Pakiety dla języka polskiego ---
\usepackage[utf8]{inputenc}
\usepackage[T1]{fontenc}
\usepackage[polish]{babel}
\usepackage{tabularx} 
\usepackage{booktabs} % Dla ładniejszych linii w tabelach
\usepackage{listings} % Dla bloków kodu
% --- Marginesy i formatowanie ---
\usepackage[margin=2.5cm]{geometry}
\usepackage{enumitem} % Lepsza kontrola nad listami

\title{Dokumentacja Funkcjonalna Projektu:\\ Generowanie Grafów Planarnych}
\author{Aleksandra Rybałko, Nella Karczewicz}
\date{\today}

\begin{document}

\maketitle

\section{Cel projektu}

\subsection{Cel główny}
Celem projektu jest implementacja programu w języku C, który automatycznie wyznacza współrzędne wierzchołków grafu planarnego w celu stworzenia jego przejrzystej i estetycznej wizualizacji. Aplikacja przyjmuje opis krawędzi grafu, a następnie, wykorzystując algorytmy \textbf{Fruchterman-Reingolda} oraz \textbf{Tutte’go}, oblicza pozycje $x$ i $y$ dla każdego węzła, dążąc do minimalizacji przecięć krawędzi. Następnie zapisuje wynik działania programu jako plik tekstowy lub binarny.

\subsection{Podcele}
Cele szczegółowe projektu obejmują następujące zadania:
\begin{itemize}[label=---]
    \item Wczytanie grafu z pliku tekstowego.
    \item Wyznaczanie współrzędnych węzłów dla przystępnej i czytelnej wizualizacji grafu planarnego przedstawionego w postaci listy krawędzi poprzez implementację dwóch algorytmów: Fruchterman-Reingolda oraz Tutte.
    \item Zapisanie wyniku działania programu jako pliku tekstowego lub binarnego (według wyboru użytkownika) zawierającego listę współrzędnych wierzchołków.
\end{itemize}

\subsection{Wymagania techniczne}
Program został zaprojektowany zgodnie z poniższymi wymaganiami technicznymi:
\begin{itemize}
    \item Kod musi być zgodny z dobrymi praktykami programowania w języku C oraz powinien być elastyczny i przejrzysty.
    \item Program powinien kompilować się przy użyciu kompilatora \texttt{GCC}.
    \item Aplikacja w języku C jest sterowana jedynie z linii poleceń, przez argumenty i opcje.
    \item Program wczytuje dane w formacie: \\
    \texttt{<nazwa\_krawędzi> <wierzchołek\_A> <wierzchołek\_B> <waga\_krawędzi>}
\end{itemize}

\paragraph{Przykład wejścia:}
\begin{verbatim}
AB;1;2;1
BC;2;3;1
CD;3;4;1
DB;4;2;1.407
\end{verbatim}

\paragraph{Przykład wyjścia:}
Wynikiem działania jest plik (\texttt{txt} lub \texttt{bin}) o strukturze: \\
\texttt{<wierzchołek> <współrzędna\_x> <współrzędna\_y>}
\begin{verbatim}
1;0.0;0.0
2;1.0;0.0
3;1.0;1.0
4;0.0;1.0
\end{verbatim}

\section{Ogólne funkcje programu}
Aplikacja odpowiada za wykonywanie następujących zadań:

\begin{enumerate}
    \item \textbf{Wczytywanie struktury grafu:} Program odczytuje dane wejściowe jako pliki tekstowe zawierające listę krawędzi wraz z ich wagami i buduje wewnętrzną reprezentację powiązań między węzłami.
    \item \textbf{Automatyczne wyznaczanie współrzędnych:} Główna funkcja programu polegająca na zamianie listy połączeń w konkretne punkty na płaszczyźnie. Program dąży do optymalnego rozłożenia wierzchołków, minimalizując ich nakładanie się na siebie.
    \item \textbf{Wybór algorytmu:} Użytkownik ma możliwość wskazania, który algorytm (Fruchterman-Reingolda lub Tutte) ma zostać użyty do obliczeń poprzez parametry wywołania.
    \item \textbf{Zarządzanie formatami danych:} Konwersja danych wejściowych na ustandaryzowany format wyjściowy, umożliwiający późniejszą prezentację wyników w aplikacjach graficznych (np. w module Java).
\end{enumerate}
\section{Specyfikacja flag sterujących (CLI)}

Program sterowany jest za pomocą argumentów przekazywanych w wierszu poleceń. Poniższa tabela przedstawia dopuszczalne flagi oraz zakresy ich parametrów:

\begin{table}[h]
    \centering
    \small
    \begin{tabularx}{\textwidth}{@{} l l l X @{}}
        \toprule
        \textbf{Flaga} & \textbf{Argument} & \textbf{Wartość} & \textbf{Opis funkcjonalny} \\
        \midrule
        \texttt{-i} & \texttt{<path>} & Ścieżka & Plik wejściowy. Lokalizacja pliku tekstowego z listą krawędzi. \\
        \texttt{-o} & \texttt{<path>} & Ścieżka & Plik wyjściowy. Nazwa i lokalizacja pliku wynikowego. \\
        \texttt{-alg} & \texttt{<id>} & 0 lub 1 & Wybór algorytmu. Pozwala na zmianę metody obliczeniowej. \\
        \texttt{-f} & \texttt{<format>} & txt lub bin & Format zapisu. \texttt{txt}: tekstowy, \texttt{bin}: binarny. \\
        \texttt{-n} & \texttt{<int>} & > 0 & Liczba iteracji. Precyzja procesu optymalizacji (domyślnie 500). \\
        \texttt{-h} & brak & brak & Pomoc. Wyświetla instrukcję obsługi w konsoli. \\
        \bottomrule
    \end{tabularx}
    \caption{Zestawienie parametrów wiersza poleceń.}
\end{table}

\subsection{Przykłady wywołania programu}

Obliczenia standardowe (format tekstowy):
\begin{verbatim}
./graph_layout -i graf.txt -o wynik.txt -alg 0 -n 1000
\end{verbatim}

Obliczenia zoptymalizowane (format binarny):
\begin{verbatim}
./graph_layout -i graf.txt -o punkty.bin -alg 1 -f bin
\end{verbatim}

\subsection{Walidacja argumentów i danych wejściowych}

Program przed rozpoczęciem fazy obliczeniowej przeprowadza weryfikację poprawności:

\begin{itemize}[label=\textbullet]
    \item \textbf{Kontrola obecności plików:} Program sprawdza, czy plik wejściowy podany po fladze \texttt{-i} fizycznie istnieje na dysku.
    \item \textbf{Weryfikacja typów:} W przypadku podania wartości nieliczbowej dla parametru \texttt{-n} lub \texttt{-alg}, program przerywa działanie, wskazując błędny argument.
    \item \textbf{Spójność danych:} Program weryfikuje, czy każda linia pliku wejściowego zawiera wymagane cztery kolumny danych.
\end{itemize}
\section{Specyfikacja formatów danych}

\subsection{Format wejściowy (Import)}

Program obsługuje import danych w formacie tekstowym opartym na strukturze \textbf{CSV/SSV}. Każda linia pliku wejściowego definiuje pojedynczą krawędź grafu i musi składać się z czterech kolumn:
\begin{center}
    \texttt{Nazwa\_krawędzi; Wierzchołek\_A; Wierzchołek\_B; Waga}
\end{center}

\paragraph{Szczegółowa walidacja pól i ograniczenia:}
W celu zapewnienia stabilności fazy obliczeniowej, nałożono następujące restrykcje na dane wejściowe:

\begin{itemize}[label=\textbullet]
    \item \textbf{Nazwa krawędzi (string):} 
        \begin{itemize}[label=---]
            \item Maksymalna długość: \textbf{10 znaków} (optymalizacja pod limit 1000 krawędzi).
            \item Dopuszczalne znaki: wyłącznie małe litery alfabetu łacińskiego (\texttt{a-z}) oraz cyfry.
            \item Zabronione: spacje, wielkie litery, znaki specjalne oraz polskie znaki diakrytyczne (\texttt{ą, ę, ł} itd.).
            \item Format: ciągły ciąg znaków bez przerw.
        \end{itemize}
    \item \textbf{Identyfikatory wierzchołków (int):} 
        \begin{itemize}[label=---]
            \item Typ: liczba całkowita dodatnia.
            \item Maksymalna długość: \textbf{4 cyfry} (zakres do 9999, przy limicie 1000 unikalnych węzłów).
            \item Zabronione: litery oraz znaki specjalne.
        \end{itemize}
    \item \textbf{Waga krawędzi (double):} 
        \begin{itemize}[label=---]
            \item Format: liczba dziesiętna z kropką jako separatorem.
            \item Maksymalna długość zapisu: \textbf{4 znaki} (np. \texttt{1.00}).
            \item Autokorekta: program automatycznie zamienia przecinek na kropkę przed procesem parsowania.
        \end{itemize}
\end{itemize}

\paragraph{Zasady interpretacji separatorów i liczb:}
Program stosuje elastyczne podejście do formatowania technicznego:
\begin{itemize}[label=\textbullet]
    \item \textbf{Separatory danych:} Jako rozdzielacz kolumn akceptowany jest średnik (\texttt{;}) lub przecinek (\texttt{,}).
    \item \textbf{Białe znaki:} Dowolna liczba spacji i tabulatorów między kolumnami jest ignorowana.
    \item \textbf{Standard EOL:} Obsługiwany jest standard \texttt{LF} (\texttt{\textbackslash n}) typowy dla systemów Unix/Linux.
\end{itemize}

\textit{Przykład poprawnego rekordu:} \texttt{krawedz12;102;105;1.50}
\subsection{Format wyjściowy (Eksport)}

Program generuje współrzędne wierzchołków w wybranym przez użytkownika formacie (\texttt{txt} lub \texttt{bin}).

\subsubsection{Format tekstowy (TXT)}
Dane wyjściowe są zapisywane w formacie kolumnowym, gdzie separatorem jest pojedyncza spacja. Każda linia reprezentuje docelowe położenie jednego węzła:
\begin{center}
    \texttt{ID\_Wierzchołka Współrzędna\_X Współrzędna\_Y}
\end{center}

\textit{Przykład pliku wynikowego:}
\begin{verbatim}
1;15.50;40.25
2;10.00;5.12
3;60.15;12.80
\end{verbatim}

\subsubsection{Format binarny (BIN)}
W przypadku wyboru formatu binarnego (flaga \texttt{-f bin}), dane zapisywane są bez separatorów tekstowych w celu optymalizacji rozmiaru pliku i szybkości odczytu przez moduł wizualizacji. Struktura rekordu binarnego:
\begin{itemize}
    \item \texttt{int} (4 bajty) – ID wierzchołka.
    \item \texttt{double} (8 bajtów) – Współrzędna X.
    \item \texttt{double} (8 bajtów) – Współrzędna Y.
\end{itemize}

\subsection{Obsługa formatów niestandardowych}
Program został wyposażony w mechanizm \textbf{ignorowania pustych linii} oraz linii zaczynających się od znaku hash (\texttt{\#}), co pozwala na dodawanie komentarzy wewnątrz plików z danymi grafu bez wpływu na proces obliczeniowy.

\section{Sytuacje wyjątkowe i błędy}

Program został zaprojektowany z myślą o odporności na błędne dane wejściowe oraz nieprawidłowe parametry wywołania. W przypadku wystąpienia sytuacji krytycznej, aplikacja przerywa działanie, wypisuje komunikat o błędzie na standardowe wyjście błędów (\texttt{stderr}) i zwraca odpowiedni kod powrotu do systemu operacyjnego.

\subsection{Tabela kodów błędów}

Poniższa tabela zawiera zestawienie komunikatów, ich przyczyn oraz przypisanych im kodów wyjścia:

\begin{table}[h]
    \centering
    \small
    \begin{tabularx}{\textwidth}{@{} l X c @{}}
        \toprule
        \textbf{Komunikat błędu (stderr)} & \textbf{Znaczenie i przyczyna} & \textbf{Kod powrotu} \\
        \midrule
        \texttt{Błąd: Nieprawidłowe flagi CLI} & Brak flag -i/-o, brak wymaganych argumentów lub użyto nieznanego formatu wyjściowego. & \textbf{1} \\
        \addlinespace
        \texttt{Błąd: Krytyczny błąd odczytu} & Błąd alokacji pamięci RAM podczas wczytywania lub niskopoziomowy problem z dostępem do pliku. & \textbf{2} \\
        \addlinespace
        \texttt{Błąd: Nie można otworzyć pliku} & Plik wejściowy nie istnieje pod wskazaną ścieżką lub brak uprawnień do odczytu. & \textbf{3} \\
        \addlinespace
        \texttt{Błąd: Nieznany algorytm} & Wybrano nazwę algorytmu inną niż dopuszczalne \textit{tutte} lub \textit{fruchterman}. & \textbf{4} \\
        \addlinespace
        \texttt{Błąd: Plik pusty} & Wskazany plik wejściowy nie zawiera żadnych danych do przetworzenia. & \textbf{5} \\
        \addlinespace
        \texttt{Błąd: Graf zbyt mały} & Graf posiada niewystarczającą liczbę elementów do poprawnego działania (wymagane $V \ge 3, E \ge 2$). & \textbf{6} \\
        \addlinespace
        \texttt{Błąd: Wynik nieplanarny} & Wykryto przecięcia krawędzi po zakończeniu procesu optymalizacji układu. & \textbf{7} \\
        \addlinespace
        \texttt{Błąd: Graf niespójny} & Test spójności DFS nie został zaliczony; graf składa się z niepołączonych składowych. & \textbf{8} \\
        \addlinespace
        \texttt{Błąd: Format wyjściowy} & Przekazano nieprawidłowy format zapisu danych (wymagany \textit{txt} lub \textit{bin}). & \textbf{9} \\
        \addlinespace
        \texttt{Błąd: Struktura danych} & Wykryto błędy w linii: pętle własne, niedozwolone znaki diakrytyczne lub zły format kolumn. & \textbf{10} \\
        \bottomrule
    \end{tabularx}
    \caption{Specyfikacja obsługi błędów i kodów wyjścia.}
\end{table}

\subsection{Sposób raportowania}

Wszystkie błędy są raportowane w ustandaryzowany sposób, co ułatwia automatyzację pracy z programem (np. przy użyciu skryptów Bash lub modułu nadrzędnego w Javie). Przykładowy format raportowania:

\begin{quote}
    \texttt{[FATAL ERROR]: Błąd: Nieprawidłowy format danych w linii 42. Oczekiwano wartości numerycznej.}
\end{quote}

\subsection{Zabezpieczenia algorytmiczne}
Program implementuje mechanizm „bezpiecznika” (\textit{watchdog}), który przerywa obliczenia, jeśli po zadanej liczbie iteracji siły w modelu sprężynowym nie dążą do równowagi (zapobieganie pętlom nieskończonym przy specyficznych układach krawędzi).



\end{document}
